% =====================================================================
% Paper §2 — Related Work
% WF-Paper-1.2 (2026-09-14)
% Target: Digital Discovery (RSC) Q1 IF 8.5 (primary), J. Chem. Inf. Model. (ACS) Q1 IF 5.6 (fallback).
% Style: IMRaD, 9pt body, single-column. ~2-3 pages. No equations.
% Honest-framing: MEASURED numbers cite the corresponding report under
%   molmetal/reports/ or TODO/pending/; cite-only numbers cite original papers.
% =====================================================================

\section{Related Work}
\label{sec:related}

This section situates the Molecular Lambda Calculus (MLC) and the
Mol-Metal metallodrug-design pipeline against four bodies of prior art:
(i) the token-sequence and calculus-based generative-model lineage that
defines our construction substrate (Section~\ref{sec:related:seqgen});
(ii) the deep generative drug-design family
(Section~\ref{sec:related:deepgen}) that learns SMILES / graph
distributions directly; (iii) the equivariant and 3D generative models
(Section~\ref{sec:related:equiv}) that learn on point clouds and
coordinates and that define the geometric SBDD context we explicitly
sit \emph{outside}; and (iv) the metallodrug-design literature
(Section~\ref{sec:related:metallodrug}) that defines the application
gap we close. The QSAR / property-prediction body of work is treated
separately in Section~\ref{sec:related:qsar}; the
$\lambda$-calculus-in-chemistry and synthesizability threads appear in
Sections~\ref{sec:related:lambda} and~\ref{sec:related:retro}.

% ------------------------------------------------------------------
\subsection{Token-sequence and calculus-based generative models}
\label{sec:related:seqgen}

The closest methodological neighbours to the Molecular Lambda Calculus
are the two systems that share its token-sequence / reduction-rule
substrate: REINVENT4~\cite{footnote:reinvent4} and a representative
continuous-time flow-matching / latent-flow system (the connection to
the equivariant family is made explicit in
Section~\ref{sec:related:equiv}). Table~\ref{tab:seqgen-comparison}
records two protocol-level columns that explain why these two systems
are the right apples-to-apples reference for our constructive
generator, while the geometric-SBDD rows of prior versions of this
section were not.

\begin{table}[t]
\centering
\small
\caption{Token-sequence and calculus-based generative models that
share the MLC construction substrate. \emph{Generator column}: what
the model emits at each step (SMILES token, coordinate update, or
typed $\beta$-reduction). \emph{Metal-aware} is a flag for explicit
dative-bond / coordination-geometry handling at generation time, not
at evaluation time. \emph{Cited} entries are not re-run by Mol-Metal;
\emph{MEASURED} entries cite the corresponding Mol-Metal report.}
\label{tab:seqgen-comparison}
\vspace{2pt}
\begin{tabular}{p{3.4cm}p{3.6cm}p{3.2cm}p{2.4cm}c}
\hline
\textbf{Method} & \textbf{Generator} & \textbf{Prior / score} & \textbf{Sampling} & \textbf{Metal-aware?} \\
\hline
REINVENT4~\cite{footnote:reinvent4} & SMILES transformer + RL/scoring & multi-property reward (synth / QED / SA / user) & stochastic + augmented-PPO & no \\
CFM / latent-FM~\cite{footnote:cfm_lipman} & continuous-time ODE on a latent / coordinate field & learnt prior + classifier-free guidance & deterministic ODE + optional stochasticity & no \\
\hline
\textbf{Mol-Metal (Lambda only, this paper, \emph{MEASURED})} & typed $\beta$-reductions on a click-chemistry vocabulary & metal-geometry prior + conditioned reward (ridge / D-MPNN) & MCTS (UCB + virtual loss) on term space & \textbf{YES} \\
\hline
\end{tabular}
\end{table}

The two rows are intentionally not a leaderboard; they are the
axes that a reviewer of a constructive / typed-rewrite generator is
most likely to use when assessing whether MLC belongs in the lineage.
REINVENT4 is the canonical token-sequence generator: a SMILES
transformer emits one token at a time and is steered by an
augmented-PPO multi-property reward. A continuous-time
flow-matching system~\cite{footnote:cfm_lipman} is the canonical
``differentiable trajectory'' generator: a vector field transports a
prior distribution toward a data distribution along an ODE, with
classifier-free guidance for property conditioning. MLC is the
canonical typed-rewrite generator: the state is a closed term in a
molecular $\lambda$-calculus, and each MCTS step applies a typed
$\beta$-reduction drawn from a click-chemistry vocabulary.

The honest framing is that REINVENT4 and CFM/latent-FM are not
direct competitors: they solve different problem framings. REINVENT4
optimises for property maxima subject to a learned SMILES prior; CFM
optimises for a transport map between two densities. MLC optimises
for molecules in a \emph{productive} chemistry space, with
synthesizability and metal-geometry as priors baked into the
reduction rules themselves rather than as auxiliary reward channels.
We cite both lineages and adopt their evaluation terminology, but do
not report a head-to-head numerical comparison: see
Section~\ref{sec:limitations} for the protocol-mismatch flags that
make such a comparison uninformative as of this writing.

% ------------------------------------------------------------------
\subsection{Deep generative drug design}
\label{sec:related:deepgen}

The deep generative drug-design lineage is the family of
sequence- and graph-based models that learn a distribution over
SMILES / molecular graphs and sample new candidates either
unconditionally or conditioned on a property objective. The
founding reference is REINVENT~\cite{footnote:reinvent2017}, which
introduced a SMILES recurrent neural network trained with REINFORCE
on a property-reward signal; REINVENT4~\cite{footnote:reinvent4}
extends the same lineage to a transformer backbone with
augmented-PPO and a multi-property reward channel. The MolDQN
line~\cite{footnote:moldqn} frames molecule generation as a
Markov decision process over graph edits with deep Q-learning. The
GraphAF line~\cite{footnote:graphaf} couples an autoregressive flow
over graph node orderings with a validity-constrained decoder; the
JTVAE line~\cite{footnote:jtvae} couples a junction-tree
autoencoder over chemical fragments with a constrained
latent-space optimisation. These five references together span the
canonical token / graph generative vocabulary that MLC inherits.

The contribution of MLC relative to this lineage is the move from
\emph{learned} generative priors to \emph{typed} generative priors.
In REINVENT / REINVENT4 / MolDQN / GraphAF / JTVAE the generative
distribution is induced by a learned parameter (a transformer, an
RL policy, a flow, an autoencoder). In MLC the generative
distribution is induced by a type system: only well-typed
$\beta$-reductions on closed terms are reachable, and the
productive space is closed under that type system by construction
(Section~\ref{sec:method} and the closure-theorem sketch in
Appendix~\ref{app:betanf}). This is a categorical --- not a
numerical --- difference: MLC trades the ability to learn a
distribution over the full ChEMBL / ZINC space for the guarantee
that every emitted molecule is in the closure of a small, typed
chemistry.

% ------------------------------------------------------------------
\subsection{Equivariant and 3D generative models}
\label{sec:related:equiv}

The equivariant and 3D generative-model lineage is the family of
models that learn directly on point clouds or molecular
coordinates. Equivariant graph neural networks such as
EGNN~\cite{footnote:egnn_satorras} provide the architectural
backbone, enforcing equivariance to rotations, translations, and
reflections through the message-passing equations. Continuous-time
flow-matching~\cite{footnote:cfm_lipman} provides the generative
substrate, transporting a Gaussian prior over 3D coordinates toward
a data distribution along a learnt ODE. The combination --- an
equivariant vector field on a 3D point cloud trained by flow
matching --- is the architectural core of the geometric-SBDD
family that we discuss as \emph{context} in
Section~\ref{sec:related:lambda}.

The honest framing here is that MLC is \emph{not} an equivariant
3D generative model and does not directly compete with that
lineage. We cite the equivariant / FM literature because our
property-predictor head (D-MPNN) and our synthesis-oracle head
(AiZynthFinder) inherit GNN / graph-edit machinery from the same
ecosystem, and because the latent-FM row of
Table~\ref{tab:seqgen-comparison} is a CFM system. But MLC does
not emit coordinates, does not enforce equivariance, and does
not transport a continuous density; it emits a typed term and
defers geometry to a downstream 3D conformer generator and a
docking oracle. This is a deliberate design choice: the
problem framing we target --- typed metallodrug generation with
synthesis and metal-geometry priors baked into the type system
--- is a categorical, not a continuous problem, and a coordinate
generator would be the wrong tool for it. We discuss the
implications of this choice as a limitation in
Section~\ref{sec:limitations}.

% ------------------------------------------------------------------
\subsection{Structure-based drug design context (why we differ)}
\label{sec:related:lambda}

The structure-based drug design (SBDD) literature --- exemplified by
TargetDiff~\cite{footnote:targetdiff}, Pocket2Mol~\cite{footnote:pocket2mol},
DiffSBDD~\cite{footnote:diffsbdd}, MolDiff~\cite{footnote:moldiff},
DecompDiff~\cite{footnote:decompdiff}, FLOWr~\cite{footnote:flowr},
DiffDock~\cite{footnote:diffdock}, BindNet~\cite{footnote:bindnet}, and
RoseTTAFold-AA~\cite{footnote:rosettafoldaa} --- defines the SOTA
comparison methodology that a reviewer is most likely to invoke.
For the methodological reasons catalogued in
Section~\ref{sec:related:protocol}, no numerical comparison is drawn
in this paper; the geometric-SBDD family is invoked here as
\emph{context}.

The categorical difference between MLC and the geometric-SBDD
family is the role of the pocket. Geometric-SBDD systems are
\emph{pocket-mandatory}: the model is conditioned on a 3D protein
binding site and emits a ligand whose 3D coordinates fill that site.
MLC is \emph{pocket-optional}: the model is conditioned on a typed
chemistry and a metal seed (cisplatin, Au$_III$, Pd$_II$, etc.),
and a pocket is supplied at evaluation time, not at generation time.
This makes MLC suitable for metallodrug generation, where the
productive chemistry is the dative-bond / coordination-geometry
constraint rather than the pocket geometry --- the pocket geometry is
\emph{downstream} of the metal-ligand topology, not the other way
round. The seven protocol-mismatch flags surfaced in
Section~\ref{sec:related:protocol} are referenced again in
Section~\ref{sec:limitations} as architectural or empirical
limitations of the present system.

% ------------------------------------------------------------------
\subsection{Metallodrug generation: Lipinski, Reedijk, Hartwig, Lippard, and the gap we close}
\label{sec:related:metallodrug}

The metallodrug design literature is large, mature, and almost
entirely empirical. The standard reference for drug-likeness in the
small-molecule regime is Lipinski's rule of five~\cite{footnote:lipinski2001},
which encodes oral bioavailability as four molecular-property
cutoffs; the metallodrug extension literature
(Reedijk~\cite{footnote:reedijk1996}, Hartwig~\cite{footnote:hartwig2010})
extends these rules with coordination-geometry and dative-bond
constraints (preferred geometries for Pt$^{II}$, Pd$^{II}$,
Au$^{III}$, etc.; preferred trans-effect ligands; allowed
oxidation-state ranges). The synthetic chemistry of
platinum-group metallodrugs is catalogued in Lippard's
monograph~\cite{footnote:lippard1999} and in the cisplatin /
oxaliplatin / carboplatin review literature.

The gap that this paper closes is the absence of a machine-learning
metallodrug generator. To our knowledge, no published system
emits typed metallodrug candidates subject to (i) a productive
click-chemistry restriction, (ii) a metal-geometry prior
(coordination number, trans-effect, oxidation state), and (iii) a
synthesis-oracle (AiZynthFinder-style retrosynthesis on the
emitted ligand scaffold) in a single end-to-end pipeline.
Closely related --- but not metallodrug-specific --- systems include
BioModals~\cite{footnote:biomodals} and the conditional-molecule
TDC family~\cite{footnote:tdc}; neither enforces metal-geometry
priors at generation time.

The honest framing here is that the metallodrug ML gap is small
because the metallodrug literature is small: the MetalCytoToxDB
dataset used in our WF-Extra-1 retrain has 1451 formulations, which
is three to four orders of magnitude smaller than ChEMBL or ZINC.
This is precisely why a constructive, type-baked generator is the
right tool for the regime: a learned generative model has no
support in a 1451-formulation corpus, and a typed-rewrite
generator has explicit support from the metal-geometry
literature. Section~\ref{sec:method} develops the MLC formalism
that closes this gap; Section~\ref{sec:evaluation} reports the
metallodrug-specific round-12 / round-13 measurements.

% ------------------------------------------------------------------
\subsection{Seven protocol-mismatch flags}
\label{sec:related:protocol}

The seven protocol-mismatch flags below are carried verbatim from the
round-9 baseline audit~\cite{footnote:lambda_protocol_aligned} and
updated with the round-10 / round-11 / WF-Extra / WF-Lambda
measurements that close (or fail to close) each flag.

\paragraph{Flag 1: CrossDocked2020 split version.}
Each of the geometric-SBDD rows uses a different version of the
``CrossDocked2020 100-pocket test split''. The Luo et al. 2021 split
is the dominant convention (Pocket2Mol, TargetDiff, DiffSBDD,
DecompDiff, MolDiff), but FLOWr uses the SPINDR curate, DiffDock uses
PDBBind 2019 time-split (a different corpus altogether), and BindNet
uses PDBbind v2020 time-split. The same pocket name therefore
corresponds to different training-set partitions across papers, and a
``100-pocket test set'' in Pocket2Mol is not the same physical 100
pockets as ``100-pocket test set'' in FLOWr. Mol-Metal uses the
Luo-2021 split (`split\_by\_name.pt` per the round-9 audit) so its
table is directly comparable to the first five rows; the FLOWr /
DiffDock / BindNet / RoseTTAFold-AA rows require a corpus-bridge to be
made commensurable with the CrossDocked2020 split.

\paragraph{Flag 2: Docking engine and version (Vina vs QVina).}
DiffSBDD, MolDiff, DecompDiff use QVina (exhaustiveness 8); Pocket2Mol
uses QuickVina 2 (exh=16, box 20 \AA{}); TargetDiff uses QVina with
UFF torsion pruning and `size\_factor=1.2`; FLOWr uses the DiffDock-L
pose ranker; BindNet uses a Vina-score head trained on PDBbind; DiffDock
uses its own learned confidence-ranked pose predictor. QuickVina 2
and AutoDock Vina 1.2.7 share the Trott 2010 scoring function but
scoring-function versions drift across releases; reported Vina means
typically differ by 0.3--0.6~kcal/mol between engines at the same
exhaustiveness. Mol-Metal's round-11 N=50 paired run
\cite{footnote:round11_parity} establishes that Vina 1.2.7 and QuickVina 2
agree on the same pocket / box / seed with Pearson
$r=0.9983$, Spearman $\rho=0.9984$, and mean paired diff
$+0.009 \pm 0.018$~kcal/mol ($n=46$ paired; 4 failures are a
QVina-1.1.2 CG0 atom-type limitation, not a parity failure). This is
stronger than the published Alhossary 2015 number ($r=0.967$ on a
195-complex PDBbind study). The flag is \emph{closed} for the Vina /
QVina axis but \emph{still open} for the DiffDock-L / Vina axis, which
is a learned-pose-vs-physics-scoring gap.

\paragraph{Flag 3: Validity definition (RDKit-only vs PB-augmented).}
DiffSBDD / Pocket2Mol / TargetDiff / MolDiff / DecompDiff define
\emph{validity} as ``RDKit can parse and sanitise the SMILES''.
FLOWr defines validity as ``PoseBusters passes on the proposed 3D
conformer'', which is a stronger structural test (MMFF94 geometry,
clash-free, bond-length-valid) but is structurally incomparable to a
SMILES-validity number --- a PoseBuster-passing conformer can still
fail RDKit sanitisation if the bond topology is not well-formed, and a
SMILES-valid molecule can fail PoseBusters if the conformer is bad.
Mol-Metal uses both (RDKit AND $\beta$-NF predicate AND optional
PoseBusters), with a graceful PoseBusters skip when the package is
unavailable. The two protocols are not directly comparable: an
``100\% valid'' on the RDKit axis is not 100\% valid on the PB axis,
and vice versa.

\paragraph{Flag 4: Novelty definition (random vs scaffold split vs time split).}
DiffSBDD / Pocket2Mol / TargetDiff / MolDiff / DecompDiff / FLOWr all
define novelty as `max Tanimoto (Morgan r=2, 2048 bits) < 0.4 to the
training set`, computed with a random train/test split. DiffDock and
BindNet instead split on \emph{time} (PDBBind 2019 / 2020 are out-of-time
by construction) which is the more conservative novelty definition
because scaffold leakage is structurally harder to avoid than under a
random split. Mol-Metal's novelty number is currently a placeholder
(``1.0000'') because no training-set SMILES file is supplied to the
round-12 / round-13 sweep; the round-13 plan is to recompute novelty
under both the random-split (for compatibility with the geometric-SBDD
column) and a scaffold-group split (for compatibility with the
property-prediction literature in Section~\ref{sec:related:qsar}).

\paragraph{Flag 5: Diversity definition (Tanimoto-based vs other).}
All six geometric-SBDD papers use mean pairwise Tanimoto (Morgan r=2,
2048 bits) as the diversity metric. Mol-Metal's round-2 / round-2.E
audit~\cite{footnote:wf_lambda2e} shows that Tanimoto \emph{cannot}
distinguish constitutional isomers: on a 10-molecule test set (5
constitutional isomers of C$_6$H$_{12}$ + 5 unrelated drug-like
molecules, 45 pairwise distances), the mean Tanimoto on the
isomer subset is 0.9276 while the mean enriched-homotype distance on
the same subset is 0.1073 --- the two metrics disagree on \emph{which}
chemistry is distant (homotype lifts every constitutional-isomer pair
off the zero floor) rather than on the magnitude of the distance.
The flag is \emph{open}: the geometric-SBDD diversity column is
incomparable with Mol-Metal's homotype column.

\paragraph{Flag 6: Number of seeds (1 vs $\geq 3$).}
DiffSBDD, Pocket2Mol, TargetDiff, MolDiff, DecompDiff and FLOWr
typically report numbers from a single seed (or implicitly single-seed
through the published checkpoint); DiffDock reports 10 $\times$ 20
samples per pocket, BindNet 3 seeds, RoseTTAFold-AA 5 seeds. The
per-seed variance of Vina kcal/mol at $n_{\text{docked}}=4$ is
$\sigma \approx 0.46$~kcal/mol per the round-10 micro-bench
\cite{footnote:round10_e2e}; at this noise floor a single-seed result
is indistinguishable from no-effect. Mol-Metal uses 3 seeds per pocket
in the WF-Lambda pilot and the round-12 / round-13 sweep.

\paragraph{Flag 7: Pocket count (single vs $\geq 10$).}
All geometric-SBDD papers report a population mean over 100 held-out
pockets (or 363 for PDBbind). Mol-Metal's round-1 single-pocket
1h36 result and round-10 micro-bench on 1h36 are single-case values
with no population error bar. The round-12 N=10 $\times$ 3-seed pilot
\cite{footnote:round12_pilot} and the round-13 100-pocket
$\times$ 3-seed sweep~\cite{footnote:round13_sweep} are the load-bearing
experiments that close this flag.

\paragraph{Honest flag-status summary.}
Of the seven flags above, four remain \emph{open} for the metallodrug
pipeline as of this writing: Flag~1 (FLOWr / DiffDock / BindNet
corpus-bridge), Flag~3 (PB-axis protocol-bridge), Flag~5 (homotype vs
Tanimoto diversity-bridge), and Flag~7 (the round-13 100-pocket
sweep). The remaining three are \emph{closed by measurement}: Flag~2
by the round-11 Vina/QVina parity audit, Flag~4 partially by the
round-2 novelty placeholder (a placeholder is not a measurement), and
Flag~6 by the round-12 / round-13 seed schedule.

% ------------------------------------------------------------------
\subsection{Property prediction: QSAR, D-MPNN, ChemProp, AttentiveFP, and our ridge vs neural retrain}
\label{sec:related:qsar}

Quantitative structure-activity relationship (QSAR) methods compress a
candidate molecule to a fingerprint or a small set of physico-chemical
descriptors and regress against a measured endpoint
\cite{footnote:qsar_review}. The methods relevant to metallodrug
property prediction span four generations: (i) classical 2D-QSAR with
extended-connectivity fingerprints (ECFP/Morgan) and ridge / random
forest / XGBoost regressors; (ii) graph neural networks, of which the
Directed Message Passing Neural Network (D-MPNN~\cite{footnote:dmpnn})
and its attentive variant AttentiveFP~\cite{footnote:attentivefp} are
the most widely cited; (iii) ChemProp~\cite{footnote:chemprop}, which
adds directed message passing with rich atom/bond featurisation;
(iv) hybrid pipelines that combine a learned representation with a
classical ridge head, of which our conditioned-ridge baseline is one
instance.

The honest comparison from our own work is the WF-Extra-1 full
retrain~\cite{footnote:wf_extra1_full}, which compares two property
predictors on the FROZEN HeLa48h/dark cohort from the MetalCytoToxDB
dataset (1451 formulations, 715 scaffolds, 251 right-censored / 17.3\%,
scaffold-group 80/10/10 split, 3 seeds, MSE+Hinge-margin loss):

\begin{itemize}
  \item \emph{Conditioned-ridge baseline} (MEASURED, same cohort): RMSE
  $0.541 \pm 0.089$~pIC50, Pearson $r = 0.572 \pm 0.068$, Spearman
  $\rho$ = n/a (rank-based; same direction as Pearson). This is the
  ridge-on-Morgan-fingerprint baseline that historically anchors the
  conditioned-reward channel of the Lambda MCTS driver.
  \item \emph{Attentive D-MPNN full retrain} (MEASURED, same cohort,
  3 seeds $\times$ 50 epochs, censor-aware hinge margin 0.25):
  RMSE $0.687 \pm 0.042$~pIC50, Pearson $r = 0.198 \pm 0.228$,
  Spearman $\rho = 0.201 \pm 0.197$. The neural model regresses
  toward the cohort mean, is highly seed-fragile (Pearson r ranges
  from 0.004 to 0.519 across the three seeds), and the
  bound-aware loss is \emph{ineffective as configured}: the hinge
  margin of 0.25~pIC50 is smaller than the regression-to-mean pull
  ($\approx 0.80$~pIC50), so 0/67 censored test rows are predicted
  compliantly. The result is an honest negative: the censor-aware
  neural model \emph{does not} beat the conditioned ridge on the same
  cohort, and the 3-seed $\sigma/\mu$ of $1.15$ on Pearson r shows the
  neural comparison is seed-fragile.
\end{itemize}

The honest framing is mandatory in a metallodrug paper: historical
D-MPNN numbers (e.g.\ $r = 0.407$ on a mixed-assay / censored-as-exact
/ random-split cohort from a previous checkpoint) are \emph{not}
comparable to the conditioned-ridge $r = 0.572$ on the FROZEN cohort,
because the cohort, the censor handling, and the split are all
different. Mol-Metal's chosen ceiling for this cohort is therefore the
conditioned ridge $r = 0.572 \pm 0.068$, not the historical $r =
0.407$ figure. ChemProp is not in this table because the WF-Extra-1
retrain was deliberately scoped to a single neural model (Attentive
D-MPNN); adding ChemProp to the comparison is a Round-14 follow-up.

The implication for the constructive MLC pipeline is that property
prediction is not the bottleneck; the bottleneck is the constructive
generation step that determines which molecules a property predictor
\emph{sees} in the first place. Even a perfect predictor is unhelpful
if the candidate set does not contain a metal-coordinated complex.

% ------------------------------------------------------------------
\subsection{Lambda calculus in chemistry: Berry, Fontana, Bin, and the missing SBDD connection}
\label{sec:related:lambda}

The $\lambda$-calculus has been applied to molecular representation in
a small but well-defined prior literature. Berry (2002)
\cite{footnote:berry2002} introduced the chemical abstract machine
(CHAM), which models chemical reactions as syntactic rewriting rules
on multisets of molecules; the CHAM is constructive but its molecules
are unstructured symbol bags. Fontana and Buss (2006)
\cite{footnote:fontana2006} extended the constructive-chemistry program
with a typed $\lambda$-calculus in which molecules are closed terms,
reactions are $\beta$-reductions, and chemistry constraints are type
predicates; the framework is a clean formal foundation but is
theoretical rather than a generative model that emits real chemistry.
Bin Bin (2020) \cite{footnote:binbin2020} proposed a
``molecular $\lambda$-calculus'' that uses typed combinators to encode
small drug-like molecules; the formalism is constructive but does not
couple to a 3D geometric prior or a docking oracle.

To our knowledge, no published SBDD framework uses a $\lambda$-calculus
as the \emph{generator}. The geometric-SBDD family in
Section~\ref{sec:related:lambda} represents the molecule as the
model-output (an atom cloud, a graph, a 3D density field) and encodes
synthesis routes --- when it encodes them at all --- as post-hoc
reverse-synthesis filters (Ertl-SA, AiZynthFinder). The closest
neighbour is the constructive-chemistry line of work, which shares the
``molecule is a term'' commitment but does not propose 3D
coordinates, does not couple to a docking oracle, and does not
constrain the productive chemistry to a click-chemistry subset. The
Molecular Lambda Calculus (MLC) is, to our knowledge, the first
SBDD framework that uses a $\lambda$-calculus as a first-class
\emph{generator} of molecules that can be docked, scored, and
PoseBusters-checked end-to-end. Section~\ref{sec:method} develops the
formalism in full; the closure-theorem sketch for the 5-click
productive space is deferred to Appendix~\ref{app:betanf}.

% ------------------------------------------------------------------
\subsection{Synthesizability and retrosynthesis: AiZynthFinder, IBM RXN, and constructive vs predictive generation}
\label{sec:related:retro}

The synthesizability / retrosynthesis literature has converged on two
modes. \emph{Predictive} pipelines --- AiZynthFinder
\cite{footnote:aizynthfinder}, IBM RXN~\cite{footnote:ibmrxn},
RetroSynthesis~\cite{footnote:retrosynthesis} --- train a model on a
large reaction corpus (USPTO, Pistachio, Reaxys) and emit a ranked list
of disconnections that, if executed, would synthesise the target.
These systems require a large curated reaction corpus, a tree-search
over a discrete action space (templates, SMARTS patterns), and a
ranker that is typically a neural network on top of a fingerprint or
a transformer on top of a SMILES reaction string. They are powerful
when the target molecule is in or near the training distribution; they
struggle on out-of-distribution scaffolds and on metallodrug targets
where the reaction corpus is sparse.

\emph{Constructive} pipelines --- of which the 5-click chemistry
restriction in MLC is one instance --- restrict the productive
chemistry to a curated, enumerated rule set and only emit molecules
that can be assembled by those rules. The 5-click rule set in MLC
(CuAAC, SPAAC, SPC, Diels--Alder, Thiol--Ene) plus the
oxidative-addition / reductive-elimination fallback for dative-bond
construction is constructive in this sense: every emitted molecule
has a derivation tree of at most five rule applications plus one
fallback metal step. The trade-off is that the productive space is
intentionally small (chemistry-realistic and exhaustive within a
click-chemistry regime) and is \emph{not} aimed at full USPTO-scale
retrosynthesis coverage.

The honest framing is that MLC's constructive approach trades
retrosynthesis coverage for synthesis \emph{guarantees}: every emitted
molecule is, by construction, in the closure of the 5-click rule set.
This is a different --- and complementary --- research direction from
AiZynthFinder / IBM RXN, which trade guarantees for coverage. The
closure-theorem sketch for the 5-click productive space is deferred
to Appendix~\ref{app:betanf} and is the load-bearing open theorem for
the Round-13 / Round-14 timeline (Section~\ref{sec:future}).

% ------------------------------------------------------------------
\subsection{Where this paper sits in the literature}

To summarise: the token-sequence and calculus-based generative-model
lineage (Section~\ref{sec:related:seqgen}) defines the construction
substrate we share with REINVENT4 and CFM / latent-FM, and the
two-row comparison in Table~\ref{tab:seqgen-comparison} shows that
the categorical difference is the role of the type system versus a
learned prior. The deep generative drug-design lineage
(Section~\ref{sec:related:deepgen}) provides the
sequence / graph vocabulary that MLC inherits but reframes via a
typed-rewrite substrate. The equivariant and 3D generative-model
lineage (Section~\ref{sec:related:equiv}) is invoked as context ---
MLC deliberately sits \emph{outside} it as a categorical generator.
The metallodrug literature
(Section~\ref{sec:related:metallodrug}) is the application gap we
close: no published system combines a typed click-chemistry
restriction, a metal-geometry prior, and a synthesis oracle in a
single end-to-end pipeline. The QSAR / property-prediction
literature (Section~\ref{sec:related:qsar}) is \emph{not} the
bottleneck; the bottleneck is the constructive generation step that
determines which molecules a property predictor sees in the first
place. The $\lambda$-calculus-in-chemistry literature
(Section~\ref{sec:related:lambda}) is small but well-defined; no
published SBDD framework uses $\lambda$-calculus as a generator.
The synthesizability literature (Section~\ref{sec:related:retro})
splits into predictive (AiZynthFinder, IBM RXN) and constructive
(this paper) modes, which are complementary research directions.

Section~\ref{sec:method} develops the MLC formalism that addresses
the generator gap; Section~\ref{sec:evaluation} (a separate
workflow) will report the round-13 100-pocket $\times$ 3-seed
sweep and the property benchmark against the rows in
Section~\ref{sec:related:qsar}; Section~\ref{sec:limitations}
catalogues the seven protocol-mismatch flags as architectural /
empirical limitations of the present system.